% GENERATED FILE. Edit canonical lessons, then run npm run generate:books.
% canonical-source-hash: fnv1a64:b5777cdb01881b1a
% canonical-lessons: HI-C432-mal, HI-C432-bhandar, HI-C432-paiket, HI-C432-lebal, HI-C432-lapetna

\chapter{Goods and Stock}
\label{ch:hi-goods-and-stock}
\hlchaptermodality{432}

\begin{chapteropening}
\textbf{By the end of this chapter:} \emph{I can talk about goods, stores and packets, read a label and get something wrapped.}

Complete the last lesson of chapter 432: I can talk about goods, stores and packets, read a label and get something wrapped.
\end{chapteropening}

\section[\texorpdfstring{\hi{माल} (māl)}{माल (māl)}]{\hi{माल} (māl) — goods; stock}

\label{lesson:HI-C432-mal}



\begin{quote}
Before the new one: say the Hindi for a sale, then the Hindi for good value.
\end{quote}

\subsection*{You'll want to know: \hi{माल}}
\textbf{\hi{माल}} — \emph{māl} — \textquotedblleft{}goods; stock\textquotedblright{}.

\begin{quote}
\hi{दुकान} \hi{में} \hi{नया} \hi{माल} \hi{आया} \hi{है।}
\end{quote}

\begin{quote}
\emph{dukān meṁ nayā māl āyā hai.}
\end{quote}

\begin{quote}
New stock has come into the shop.
\end{quote}

Masculine.

\begin{cousinweb}
\textbf{\hi{माल}} (\emph{māl}) is usually traced to Arabic \emph{māl}, \textquotedblleft{}property, wealth\textquotedblright{}, through Persian; in Hindi it is what a shop has to sell. \hi{मालगाड़ी} is a goods train.
\end{cousinweb}

\begin{grammarlens}[title={What you've built}]
One more word for getting around.
\end{grammarlens}

\subsection*{Your turn}
\begin{itemize}
\raggedright
  \item \emph{Say it:} \emph{māl}
  \item \emph{Say it:} \emph{māl}, once more
  \item \emph{Say it:} \emph{dukān meṁ nayā māl āyā hai.} — New stock has come into the shop.
  \item \emph{Recall:} say the Hindi for a sale, then the Hindi for good value, then say \emph{māl} again
\end{itemize}

\begin{tcolorbox}[breakable,colback=teal!4,colframe=teal!35!black,title={Before you move on}]
What does \hi{माल} mean? (\textquotedblleft{}Goods; stock\textquotedblright{}.) Say it once more.
\end{tcolorbox}
\section[\texorpdfstring{\hi{भंडार} (bhaṇḍār)}{भंडार (bhaṇḍār)}]{\hi{भंडार} (bhaṇḍār) — a store; a stock}

\label{lesson:HI-C432-bhandar}



\begin{quote}
Before the new one: say the Hindi for good value, then the Hindi for goods.
\end{quote}

\subsection*{You'll want to know: \hi{भंडार}}
\textbf{\hi{भंडार}} — \emph{bhaṇḍār} — \textquotedblleft{}a store; a stock\textquotedblright{}.

\begin{quote}
\hi{गाँव} \hi{के} \hi{भंडार} \hi{में} \hi{बहुत} \hi{अनाज} \hi{है।}
\end{quote}

\begin{quote}
\emph{gā̃v ke bhaṇḍār meṁ bahut anāj hai.}
\end{quote}

\begin{quote}
There is plenty of grain in the village store.
\end{quote}

Masculine.

\begin{cousinweb}
\textbf{\hi{भंडार}} (\emph{bhaṇḍār}) is a tadbhava of Sanskrit \emph{bhāṇḍāgāra}, \textquotedblleft{}a treasury\textquotedblright{}, literally a house for goods (CDIAL 9442), which lists Hindi \emph{bhā̃ḍār}, \textquotedblleft{}a storehouse\textquotedblright{}.
\end{cousinweb}

\begin{grammarlens}[title={What you've built}]
One more word for getting around.
\end{grammarlens}

\subsection*{Your turn}
\begin{itemize}
\raggedright
  \item \emph{Say it:} \emph{bhaṇḍār}
  \item \emph{Say it:} \emph{bhaṇḍār}, once more
  \item \emph{Say it:} \emph{gā̃v ke bhaṇḍār meṁ bahut anāj hai.} — There is plenty of grain in the village store.
  \item \emph{Recall:} say the Hindi for good value, then the Hindi for goods, then say \emph{bhaṇḍār} again
\end{itemize}

\begin{tcolorbox}[breakable,colback=teal!4,colframe=teal!35!black,title={Before you move on}]
What does \hi{भंडार} mean? (\textquotedblleft{}A store; a stock\textquotedblright{}.) Say it once more.
\end{tcolorbox}
\section[\texorpdfstring{\hi{पैकेट} (paikeṭ)}{पैकेट (paikeṭ)}]{\hi{पैकेट} (paikeṭ) — a packet}

\label{lesson:HI-C432-paiket}



\begin{quote}
Before the new one: say the Hindi for goods, then the Hindi for a store.
\end{quote}

\subsection*{You'll want to know: \hi{पैकेट}}
\textbf{\hi{पैकेट}} — \emph{paikeṭ} — \textquotedblleft{}a packet\textquotedblright{}.

\begin{quote}
\hi{मुझे} \hi{चाय} \hi{का} \hi{एक} \hi{पैकेट} \hi{दीजिए।}
\end{quote}

\begin{quote}
\emph{mujhe cāy kā ek paikeṭ dījie.}
\end{quote}

\begin{quote}
Give me a packet of tea, please.
\end{quote}

Masculine.

\begin{cousinweb}
\textbf{\hi{पैकेट}} (\emph{paikeṭ}) is borrowed from English \emph{packet}.
\end{cousinweb}

\begin{grammarlens}[title={What you've built}]
One more word for getting around.
\end{grammarlens}

\subsection*{Your turn}
\begin{itemize}
\raggedright
  \item \emph{Say it:} \emph{paikeṭ}
  \item \emph{Say it:} \emph{paikeṭ}, once more
  \item \emph{Say it:} \emph{mujhe cāy kā ek paikeṭ dījie.} — Give me a packet of tea, please.
  \item \emph{Recall:} say the Hindi for goods, then the Hindi for a store, then say \emph{paikeṭ} again
\end{itemize}

\begin{tcolorbox}[breakable,colback=teal!4,colframe=teal!35!black,title={Before you move on}]
What does \hi{पैकेट} mean? (\textquotedblleft{}A packet\textquotedblright{}.) Say it once more.
\end{tcolorbox}
\section[\texorpdfstring{\hi{लेबल} (lebal)}{लेबल (lebal)}]{\hi{लेबल} (lebal) — a label}

\label{lesson:HI-C432-lebal}



\begin{quote}
Before the new one: say the Hindi for a store, then the Hindi for a packet.
\end{quote}

\subsection*{You'll want to know: \hi{लेबल}}
\textbf{\hi{लेबल}} — \emph{lebal} — \textquotedblleft{}a label\textquotedblright{}.

\begin{quote}
\hi{पैकेट} \hi{पर} \hi{लेबल} \hi{पढ़िए।}
\end{quote}

\begin{quote}
\emph{paikeṭ par lebal paṛhiye.}
\end{quote}

\begin{quote}
Read the label on the packet.
\end{quote}

Masculine.

\begin{cousinweb}
\textbf{\hi{लेबल}} (\emph{lebal}) is borrowed from English \emph{label}.
\end{cousinweb}

\begin{grammarlens}[title={What you've built}]
One more word for getting around.
\end{grammarlens}

\subsection*{Your turn}
\begin{itemize}
\raggedright
  \item \emph{Say it:} \emph{lebal}
  \item \emph{Say it:} \emph{lebal}, once more
  \item \emph{Say it:} \emph{paikeṭ par lebal paṛhiye.} — Read the label on the packet.
  \item \emph{Recall:} say the Hindi for a store, then the Hindi for a packet, then say \emph{lebal} again
\end{itemize}

\begin{tcolorbox}[breakable,colback=teal!4,colframe=teal!35!black,title={Before you move on}]
What does \hi{लेबल} mean? (\textquotedblleft{}A label\textquotedblright{}.) Say it once more.
\end{tcolorbox}
\section[\texorpdfstring{\hi{लपेटना} (lapeṭnā)}{लपेटना (lapeṭnā)}]{\hi{लपेटना} (lapeṭnā) — to wrap}

\label{lesson:HI-C432-lapetna}



\begin{quote}
Before the new one: say the Hindi for a packet, then the Hindi for a label.
\end{quote}

\subsection*{You'll want to know: \hi{लपेटना}}
\textbf{\hi{लपेटना}} — \emph{lapeṭnā} — \textquotedblleft{}to wrap\textquotedblright{}.

\begin{quote}
\hi{दुकानदार} \hi{ने} \hi{किताब} \hi{काग़ज़} \hi{में} \hi{लपेटी।}
\end{quote}

\begin{quote}
\emph{dukāndār ne kitāb kāġaz meṁ lapeṭī.}
\end{quote}

\begin{quote}
The shopkeeper wrapped the book in paper.
\end{quote}

Past, with \hi{ने} (he or she …): \hi{उसने} \hi{तोहफ़ा} \hi{काग़ज़} \hi{में} \hi{लपेटा}, \emph{usne tohfā kāġaz meṁ lapeṭā}.

\begin{cousinweb}
\textbf{\hi{लपेटना}} (\emph{lapeṭnā}) goes back to \emph{*lappēṭṭ}, \textquotedblleft{}wrap\textquotedblright{}, a form CDIAL reconstructs (CDIAL 10942), which lists Hindi \emph{lapeṭnā}, \textquotedblleft{}to wrap, fold, wind\textquotedblright{}.
\end{cousinweb}

\begin{grammarlens}[title={What you've built}]
One more word for getting around.
\end{grammarlens}

\subsection*{Your turn}
\begin{itemize}
\raggedright
  \item \emph{Say it:} \emph{lapeṭnā}
  \item \emph{Say it:} \emph{lapeṭnā}, once more
  \item \emph{Say it:} \emph{dukāndār ne kitāb kāġaz meṁ lapeṭī.} — The shopkeeper wrapped the book in paper.
  \item \emph{Recall:} say the Hindi for a packet, then the Hindi for a label, then say \emph{lapeṭnā} again
\end{itemize}

\begin{tcolorbox}[breakable,colback=teal!4,colframe=teal!35!black,title={Before you move on}]
What does \hi{लपेटना} mean? (\textquotedblleft{}To wrap\textquotedblright{}.) Say it once more.
\end{tcolorbox}
